\documentclass[11pt]{article}
\usepackage[margin=1in]{geometry}
\usepackage{amsmath,amssymb,amsthm,hyperref}
\newtheorem{theorem}{Theorem}
\newcommand{\QQ}{\mathbb Q}
\title{Sparse sources span every polynomial of a complete generic binary packet}
\author{Project research note}
\date{17 November 2026}
\begin{document}
\maketitle

Let $L=2^a\ge2$, $d=2L-1$, and index rows and source bits by
$0,\ldots,L-1$. In the notation of checkpoint 009 put
\[
 C(T)=\prod_{h=-L}^{L-1}(T-h),\quad c=C'(0),\qquad
 s_x(T)=\frac{C(T)}c\sum_{y\in X}\frac1{T-(y-x)} .
\]
Here $X$ is the set of bits equal to one. A literal target assignment $Y$
gives the polynomial
\[
 F_{X,Y}=\prod_{x\in Y}s_x\prod_{x\notin Y}(1-s_x).
\]

\begin{theorem}
At every such width, the rational span of $F_{X,Y}$ with
$1\le |X|\le 2$ and unrestricted $Y$ is exactly $\QQ[T]_{\le Ld}$.
There is a polynomial-size, explicitly constructible spanning list of
target marginals with these sparse sources.
\end{theorem}

\paragraph{Reduction to $d$ coefficients at each order.}
For any row set $A$, the marginal $\prod_{x\in A}s_x$ is the sum of
$F_{X,Y}$ over $Y\supseteq A$. This includes the constant one for empty
$A$. For $|A|=k$ let
\[
 R_{A,X}(T)=\prod_{x\in A}\sum_{y\in X}\frac1{T-(y-x)} .
\]
The product $C^k R_{A,X}/c^k$ has degree at most $kd$. Modulo polynomials
of degree at most $(k-1)d$, its top $d$ coefficients are an invertible
triangular transform of the first $d$ coefficients of $T^kR_{A,X}$ at
infinity. Indeed $C$ is monic and $c\ne0$. It suffices by induction on
$k$ to show that these $d$ coefficient vectors span $\QQ^d$ at each order.

A functional on the $d$ coefficients takes the form
$R\mapsto[T^{-1}](P R)$, where $\deg P\le d+k-2$. Terms of $P$ of
degree at most $k-2$ give the zero functional. Assume the functional
vanishes on the marginals below. We show that it is zero.

\paragraph{Singletons fix the values at the grid.}
Let $G(T)=\prod_{j=1-L}^{L-1}(T-j)$, of degree $d$.
A singleton source and an interval of $k$ rows realize
$R=1/\prod_{j=h}^{h+k-1}(T-j)$ for every consecutive $k$-point
window in the grid of $G$. By partial fractions, $[T^{-1}]PR$ is the
divided difference of $P$ on this window. For $k=1$ its vanishing forces
$P=0$. For $k\ge2$ the vanishing $(k-1)$st forward differences imply
that on the grid the values agree with one polynomial of degree at most
$k-2$. Interpolate the first $k-1$ points and propagate the equality
one point at a time. Subtract that polynomial without changing the
functional. We can thus take $P=GQ$ with $\deg Q\le k-2$.
The partial-fraction formula for divided differences is also recorded
in \href{https://dlmf.nist.gov/3.3.E35}{NIST DLMF, equation 3.3.35}.

\paragraph{Double poles give equations on $Q$.}
Use $X=\{y,y+g\}$ with $g$ odd. The numerator of each row factor before
cancellation is $2(T+x-y)-g$, which is nonzero at every integer.
The poles of $R$ have order at most two since each source location
contributes at most one row at a given pole. In $GQR$, simple poles of
$R$ disappear. A double pole $t$ contributes the residue
$G'(t)Q(t)\,[ (T-t)^2R ]_{T=t}$ with nonzero multiplier of $Q(t)$.
The coefficient $[T^{-1}]GQR$ is the sum of these finite residues by
partial fraction expansion, including when $GQR$ has a polynomial part.
Consequently a window of double poles with all but one value of $Q$
already zero forces the remaining value to vanish.

For $A=\{s,\ldots,s+k-1\}$ and $g<k$, write $n=k-g$.
The double poles form the entire $n$-point interval ending at $y-s$.
When an interval of consecutive known zeros has length at least $n-1$
and lies in the union of a contiguous sliding family of $n$-point windows,
one can extend it to either end of that union using a window with a
single new point. We use this only for nonempty intervals and $n\ge2$.

\paragraph{Orders below the width.}
For $2\le k<L$ use just positions $0,\ldots,k$ for the pair tests.
Choose $s=0,1$ and $0\le y\le k-g=n$. Thus the windows have endpoints
$-1,\ldots,n$ and their union is $[-n,n]$ on the integer grid.
For even $k$, start at $g=k-1$. The single-point windows force
$Q(-1)=Q(0)=Q(1)=0$. For odd $k\ge3$, first take $g=k$, $y=0$ and
$A=\{0,\ldots,k\}\setminus\{1\}$. There is exactly one double pole,
at zero, with nonzero multiplier. Now take the consecutive rows and
$g=k-2$. Their two-point windows extend this zero through $[-2,2]$.
In both cases decrease the odd gap by two until $g=1$.
At each step the preceding interval is $[-(n-2),n-2]$, long enough
to propagate the new $n$-point windows. At the final step there are
at least $k-1$ distinct zeros. Hence $Q=0$.

\paragraph{The top order.}
Now $k=L$ is even, $A$ is fixed, and $s=0$. For $g=L-n$ odd, the
available windows end at $0,\ldots,n-1$ and cover $[-(n-1),n-1]$.
Set $v_r=|G'(r)|Q(r)$. Up to a common nonzero scalar, the residue
equations for each overlap $n$ are
\[
 \sum_{j=0}^{n-1} h_{n,j} v_{y-j}=0,\qquad 0\le y<n ,
 \quad h_{n,0}=1 ,
\]
\[
 \frac{h_{n,j+1}}{h_{n,j}}
 =\frac{(L-n+2j+2)(L-1-j)(n-1-j)}
 {(L+n-2j-2)(j+1)(L-n+1+j)} .
\]
Indeed the double-pole coefficient at $y-j$ before multiplication
by $G'$ is
\[
 \frac{\prod_{x=0}^{L-1}(2x-2j-g)}
 {(-1)^g j!(L-1-j)!(j+g)!(L-1-j-g)!}.
\]
Successive coefficients alternate in sign, as do the values of $G'$.
Their absolute-value ratios give the formula above. The weights are
positive and symmetric by reflection.

Overlap one gives $v_0=0$. This suffices at $L=2$. Write the weights
at overlap three as $(1,A,1)$. Its equations give $v_{-1}=-v_1$,
$v_2=-Av_1$ and $v_{-2}=-v_2$. For $L=4$ we have $\deg Q\le2$.
Because $|G'|$ is symmetric, $Q(0)=0$ and $Q(-1)=-Q(1)$ force $Q$
to be odd and linear. A nonzero such $Q$ gives $v_2$ and $v_1$ the
same sign, contradicting the equation.

For $L\ge8$ name the first halves of the overlap-five and overlap-seven
weight vectors $(1,B,C_5)$ and $(1,D,E,F)$. At overlap five the equations
ending at $3,4$ give
\[
 v_3=(AB-C_5+1)v_1,\qquad v_4=-Bv_3-C_5v_2-Bv_1 .
\]
Substitute into the overlap-seven equation ending at $4$. Its left-hand
side is $v_1$ times
\[
 (D-B)(AB-C_5+1)-(E-C_5-1)A+F-D-B
 =\frac{384L(L-5)(L-1)}{(L-4)(L-2)(L+1)^2(L+3)^2(L+5)}>0 .
\]
For direct verification the six weights from the ratio formula are
\begin{align*}
 A&=\frac{2(L-1)^2}{(L-2)(L+1)},&
 B&=\frac{4(L-3)(L-1)}{(L-4)(L+3)},\\
 C_5&=\frac{6(L-2)(L-1)^2}{(L-4)(L+1)(L+3)},&
 D&=\frac{6(L-5)(L-1)}{(L-6)(L+5)},\\
 E&=\frac{15(L-3)(L-2)(L-1)}{(L-6)(L+3)(L+5)},&
 F&=\frac{20(L-3)^2(L-2)(L-1)^2}{(L-6)(L-4)(L+1)(L+3)(L+5)} .
\end{align*}
It follows that $v_1=0$ and the entire interval $[-2,2]$ is zero.
The sliding windows for overlaps $5,7,\ldots,L-1$ extend this through
$[-(L-2),L-2]$. Each preceding interval has length at least $n-1$.
This forces $Q=0$ at the top order as well and completes the proof.

\paragraph{Constructive size and cost.}
There are $O(L)$ singleton windows at each order and $O(k^2)$ two-source
windows at order $k$, for $O(L^3)$ marginal descriptions over all orders.
Each description has a product of at most $L$ explicit polynomials of
degree $d$. After clearing $c^k$ their coefficients have
$O(L^2\log(L+1))$ bits by multiplying bounds on coefficient one-norms.
Generating the top coefficients and selecting a basis by exact rational
linear algebra therefore takes polynomially many bit operations in $L$.
A modular minor of full rank also certifies a layer over the rationals.
The executable uses this latter test at its listed finite widths.

A basis of literal assignments can likewise be found in polynomial time
if required. Maintain a span of chosen literals. For any marginal outside
it, split a free target bit into zero and one. The two new partial marginals
sum to their parent, so at least one lies outside the span. Choose that
child and repeat until a single literal remains, then add it. Rescan the
marginal list until it lies in the span. There are at most $Ld+1$
additions. Partial marginals remain explicit products of selected $s_x$
and $1-s_x$, with polynomial coefficient lengths. This refinement is an
algorithmic corollary. The supplied verifier works directly with marginals.

\paragraph{Consequence for precision.}
Products of $p$ independent complete generic packets span
$\QQ[T]_{\le pL(2L-1)}$ by multiplying monomial expansions. Let $M=pL$.
Any fixed exact rational linear jet readout valid on every such product
is valid on every polynomial of the degree used in checkpoint 008.
Its consecutive-grid theorem now applies to this actual generic family.
For $1\le u\le\log_2L$, $M\ge4u$, and per-node loss budgets
$\lfloor\lambda/2^{v_2(\tau)}\rfloor$, it requires
\[
 \lambda\ \ge\ \frac{ML}{8u}.
\]
For $u\ge2$ the matching order upper bound was already proved there.
This closes the smaller-rational-space route for this readout model
at growing widths. Coefficients in the span can have even denominators.
Parity recovery on the integral family, fixed C38 predicates at growing
width, and source-dependent joins are separate questions. The construction
changes the cost of certifying generic rank. The project has no improved
end-to-end ECDLP exponent from this result.

\paragraph{Reproduction.}
Run \texttt{python -B research/generic\_packet\_span/verify\_span.py}.
The checks compare literal polynomials from checkpoint 009 to marginal
coefficients, compute finite-field ranks layer by layer, and independently
evaluate the endpoint identity at integer widths. The all-width conclusion
uses the proof above.

The verified full dimensions at widths $2,4,8,16,32,64$ were
$7,29,121,497,2017,8129$, respectively. The width-64 computation
examined $17{,}553$ marginal descriptions. At width $128$ the verifier
checked only the highest layer, obtaining rank $255$ from $1{,}154$
descriptions. The literal cross-check covered $234$ small marginals.
\end{document}
